\section{INTRODUCTION}

\subsection{Contextualization And Motivation}

% missing content of vibration control [OK]
% missing fine cut vibration control > pasive mitigation []
% missing image of TLP []
% missing reference for dynamic behavior blah-blah (maybe use chakrabartha[?] 2005) [OK]

Offshore platforms operate in harsh marine environments, continuously exposed to wave action, which induces dynamic loads on their structural components -- particularly on the support legs and tendons of Tension Leg Platforms (as seen on Figure \ref{fig:TLP}). These loads propagate as mechanical vibrations throughout the structure, potentially compromising its integrity, operational performance, and service life \cite{Chakrabarti2005}.

% -----------------------------------------------
\begin{figure}[htbp] 
    \caption{Example of Tension Leg Platform.}
    \centering
    \includegraphics[scale=.7]{TLP.png}
    {\\ \footnotesize\textbf{Fonte}: Petrobrás.}
    \label{fig:TLP}
\end{figure}
% -----------------------------------------------

% To ensure the reliability and longevity of such systems, vibration mitigation has become a key research focus. Among passive control strategies, one finds 

To ensure the reliability and longevity of such systems, vibration mitigation has become a key research focus. Among passive control strategies, one finds the extensive implementation of tuned vibration absorbers, which have been rigorously investigated for deep-water applications. Considerable research efforts have demonstrated the efficacy of auxiliary damping devices --ranging from vertical caissons and Tuned Liquid Column Dampers (TLCD) to multiple Tuned Mass Damper (TMD) configurations, in suppressing resonant heave and pitch motions under stochastic wave excitation \cite{Kandasamy2016}.

% An alternative means to passive mitigation control is the introduction of intrinsic structural modifications. This shift in perspective leads to the investigation of periodic structures, also known as phononic or elastic metamaterials. Where the geometric or material design of the tendon itself acts as a passive filter for wave propagation, due the formation of stop bands. The lack of individual parts requiring maintenance turns their use advantageous in comparison to other approaches for passive vibration.

An alternative mean for passive mitigation lies in the introduction of intrinsic structural modifications. This perspective motivates the investigation of periodic structures, also called phononic or elastic metamaterials, where the periodic modulation of the tendon's geometry or material distribution acts as a passive filter for wave propagation. By inducing attenuation bands, these systems eliminate the need for complex auxiliary mechanisms. Together, these characteristics make this approach particularly advantageous for ensuring long-term robustness of TLP tendons.

% periodic structures, also known as phononic or elastic metamaterials, have gained increasing attention due to their ability to generate band gaps, \textit{i.e.}, specific frequency ranges in which wave propagation is inherently suppressed. This phenomenon arises from the structural periodicity and is strongly influenced by the geometry, boundary conditions, and material distribution along the structure.

Consequently, understanding the underlying effects and optimizing these periodic configurations can lead to increased tunability, \textit{i.e.}, control of the bandwidth and location of attenuation bands. This leads to improved vibration control in such offshore platforms, allowing for safer, more durable, and sustainable marine operations.

\subsection{Work Details}

Section 1 presents context and motivation for the development of the current study. Section 2 shows a brief history of TLP tendon modeling focused on the use of various beam theories. Then, section 3 presents the EBT model developed via analytical and FEM formulations. While section 4 presents validations of relevant routines, and the simulation of a tendon-like structure. Finally, section 5 displays the main conclusions drawn from this research.
